% ============================================================
%  Alcance del proyecto
% ============================================================
\section{Alcance del proyecto}
\label{sec:alcance}

El proyecto contempla el diseño y la construcción de un sistema integral de
ensayos de seguridad eléctrica, compuesto por un dispositivo físico y una interfaz
de software, destinado a verificar el cumplimiento de la norma IEC~60335. El
alcance abarca tanto los aspectos funcionales como no funcionales, además de las
restricciones, supuestos y riesgos asociados al desarrollo.

\subsection{Requerimientos funcionales y no funcionales}

\subsubsection{Requerimientos funcionales}
\begin{itemize}
    \item Ejecución de los ensayos de consumo, rigidez dieléctrica, continuidad de
          tierra y corrientes de fuga (activa y pasiva).
    \item Automatización de la configuración, conmutación, adquisición y
          procesamiento de datos.
    \item Generación de reportes digitales trazables y exportables.
    \item Interfaz gráfica intuitiva para la configuración y supervisión de
          ensayos.
\end{itemize}

\subsubsection{Requerimientos no funcionales}
\begin{itemize}
    \item \textbf{Confiabilidad:} resultados reproducibles con un error acotado
          respecto de los límites de la norma IEC~60335.
    \item \textbf{Seguridad:} protecciones contra sobrecargas, aislamiento
          adecuado, alarmas y comportamiento seguro ante fallas.
    \item \textbf{Usabilidad:} intervención mínima del operario y curva de
          aprendizaje baja.
    \item \textbf{Escalabilidad:} posibilidad de incorporar nuevos ensayos o
          adaptarse a modificaciones normativas.
    \item \textbf{Mantenibilidad:} arquitectura modular de hardware y software que
          facilite la reparación, actualización y extensión del sistema.
\end{itemize}

\subsection{Restricciones}
\begin{itemize}
    \item Presupuesto limitado para la construcción del prototipo y la adquisición
          de componentes.
    \item Disponibilidad restringida de equipamiento especializado para
          calibración y validación.
    \item Tiempo máximo de ejecución establecido en el cronograma del proyecto
          (Anexo~\ref{anx:cronograma}).
    \item Cumplimiento de las normativas IEC aplicables.
\end{itemize}

\subsection{Supuestos}
\begin{itemize}
    \item Se contará con un laboratorio equipado para la ejecución de ensayos de
          validación.
    \item Los componentes electrónicos necesarios estarán disponibles en el
          mercado local o mediante proveedores confiables.
    \item El personal que utilice el sistema tendrá capacitación básica en
          seguridad eléctrica y operación de equipos de ensayo.
\end{itemize}

\subsection{Riesgos y mitigación}
\begin{itemize}
    \item \textbf{Retrasos en la adquisición de componentes:} pueden comprometer el
          cronograma. Se mitigan identificando tempranamente proveedores
          alternativos.
    \item \textbf{Fallas en la calibración o en la validación del sistema:} podrían
          impedir la aceptación del proyecto. Se mitigan con pruebas piloto e
          intercomparaciones con equipos certificados.
    \item \textbf{Limitaciones presupuestarias:} aumentos de costos podrían superar
          lo estimado. Se mitigan con un margen de contingencia.
    \item \textbf{Problemas de integración entre hardware y software:} se mitigan
          con un desarrollo incremental y pruebas de integración tempranas.
\end{itemize}

\subsection{Aspectos fuera del alcance}

El proyecto se centra en el diseño, construcción y verificación de un sistema
semiautomatizado de ensayos eléctricos orientado a la norma IEC~60335. No forman
parte del alcance:
\begin{itemize}
    \item La certificación oficial del equipo ante organismos nacionales o
          internacionales (OAA, IEC, UL, entre otros), aunque el diseño se orienta
          a cumplir con dichas normativas.
    \item Ensayos no contemplados en la IEC~60335, como medición de armónicos,
          pruebas de sobrecarga o compatibilidad electromagnética.
    \item La fabricación industrial en serie del dispositivo; el trabajo se limita
          a un prototipo funcional.
    \item El mantenimiento y soporte a largo plazo del sistema.
    \item La integración con sistemas de gestión externos (ERP, MES, etc.), más
          allá del almacenamiento local y la exportación de datos.
\end{itemize}
